\section{Degenerate constructions}
\label{app:degenerate}

A statement of the form ``there is a quantum field theory on $\R^4$ with a
mass gap'' is easy to satisfy if nothing ties the theory to the Yang--Mills
action: a generalised free field whose K\"all\'en--Lehmann measure is supported
in $[\Delta^2,\infty)$ obeys the Wightman axioms and has
$\sigma(H)\subset\{0\}\cup[\Delta,\infty)$. The problem as posed by Jaffe and
Witten~\cite{JW06} excludes such answers by asking for the local gauge-invariant
fields of the Yang--Mills theory and for its short-distance behaviour. This
appendix records what happens to three natural shortcuts that stay inside the
lattice theory: the limit of the parallel transporters themselves, without
renormalisation; a theory built from two-dimensional Yang--Mills theory on
parallel planes; and a lattice theory at fixed spacing. Each gives either a
trivial theory or one that fails the axioms, and the first two statements are
theorems. None of them is used elsewhere in the paper.

Throughout, $G\subset\U(N)$ is the gauge group, $\|\cdot\|$ is the operator
norm, $\mu_L$ is the Wilson measure of the torus $\Lambda_L$ with weight
$\exp(N\beta\sum_p\ReTr Q_p)$, $\beta\ge0$, as in Section~\ref{sec:lattice},
and each link lies in six plaquettes.

\subsection{One link}

Let $\pi$ be a unitary representation of $G$ on a finite-dimensional space $V$
such that no nonzero vector is annihilated by $d\pi(\g)$; every nontrivial
irreducible representation has this property, since a vector annihilated by
$d\pi(\g)$ is fixed by the connected group $G$ and spans an invariant line.
Put
\[
D_\pi=\sup_{\|X\|\le1}\|d\pi(X)\|,\qquad
\kappa_\pi=\inf_{|v|=1}\ \sup_{\|X\|\le1}|d\pi(X)v|,
\]
the suprema over $X\in\g$. The function $v\mapsto\sup_{\|X\|\le1}|d\pi(X)v|$
is continuous and positive on the unit sphere, so $\kappa_\pi>0$, and
$\kappa_\pi\le D_\pi$. Let
\[
s_\pi=\min\Bigl\{1,\frac{\kappa_\pi}{2D_\pi^2}\Bigr\},\qquad
c_\pi=\frac{s_\pi^2\kappa_\pi^2}{16(1+e^2)}.
\]

\begin{lemma}\label{lem:onelinkdecay}
Let $\varphi\colon G\to\R$ satisfy $|\varphi(Qg)-\varphi(Q)|\le\Lambda\|g-1\|$
for all $Q,g\in G$, and let $\nu$ be the probability measure proportional to
$e^{\varphi(Q)}dQ$. Then
\[
\Bigl\|\int\pi(Q)\,d\nu(Q)\Bigr\|\le1-\frac{c_\pi}{\max(1,\Lambda)^2}.
\]
\end{lemma}

\begin{proof}
Write $M=\int\pi\,d\nu$, let $u,v$ be unit vectors and
$I=\int|\pi(Q)v-u|^2d\nu(Q)$, so that $\operatorname{Re}\ip{u}{Mv}=1-I/2$.
For $g\in G$ and every $Q$,
\[
|\pi(g)v-v|=|\pi(Qg)v-\pi(Q)v|\le|\pi(Qg)v-u|+|\pi(Q)v-u|,
\]
hence $|\pi(g)v-v|^2\le2\int|\pi(Qg)v-u|^2d\nu(Q)+2I$. Haar measure is right
invariant, so for $F\ge0$
\[
\int F(Qg)e^{\varphi(Q)}dQ=\int F(Q)e^{\varphi(Qg^{-1})}dQ
\le e^{\Lambda\|g-1\|}\int F(Q)e^{\varphi(Q)}dQ,
\]
where $\|g^{-1}-1\|=\|g-1\|$ was used. Therefore
\begin{equation}\label{eq:onelinkI}
|\pi(g)v-v|^2\le2\bigl(1+e^{\Lambda\|g-1\|}\bigr)I.
\end{equation}
Choose $X\in\g$ with $\|X\|\le1$ and $|d\pi(X)v|\ge\kappa_\pi$, put
$s=s_\pi/\max(1,\Lambda)\le1$ and $g=e^{sX}$. Then $\|g-1\|\le e^s-1\le2s$, so
$\Lambda\|g-1\|\le2$. Since $sD_\pi\le\kappa_\pi/(2D_\pi)\le\frac12$ and
$e^x-1-x\le x^2$ on $[0,1]$,
\[
|\pi(g)v-v-s\,d\pi(X)v|\le e^{sD_\pi}-1-sD_\pi\le s^2D_\pi^2,
\]
and $s\le\kappa_\pi/(2D_\pi^2)$ gives $|\pi(g)v-v|\ge s\kappa_\pi-s^2D_\pi^2\ge
s\kappa_\pi/2$. With~\eqref{eq:onelinkI}, $I\ge s^2\kappa_\pi^2/(8(1+e^2))$,
that is, $\operatorname{Re}\ip{u}{Mv}\le1-c_\pi/\max(1,\Lambda)^2$. The
supremum over $u$ and $v$ is $\|M\|$.
\end{proof}

\subsection{Straight segments}

\begin{proposition}\label{prop:segments}
Let $e_1,\dots,e_n$, $n\ge3$, be consecutive collinear links of $\Lambda_L$,
$L\ge n+1$, and $U=Q_{e_1}\cdots Q_{e_n}$. Let $\pi$ be as in
Lemma~\ref{lem:onelinkdecay}, and let $\Phi(A;Q')$ be linear in
$A\in\operatorname{End}(V)$ and bounded measurable in the links
$Q'=(Q_e)_{e\ne e_2,\dots,e_{n-1}}$, with $|\Phi(A;Q')|\le C\|A\|$. Then
\[
\bigl|\mu_L\bigl(\Phi(\pi(U);Q')\bigr)\bigr|\le
C\Bigl(1-\frac{c_\pi}{\max(1,6N^2\beta)^2}\Bigr)^{n-2}.
\]
\end{proposition}

\begin{proof}
Two distinct links lie in a common plaquette only if they are perpendicular
and share a site, or are parallel and differ by one step in a transverse
direction; collinear links never do. So every plaquette contains at most one
of the links $e_2,\dots,e_{n-1}$, and conditionally on $Q'$ these links are
independent, the link $e$ having the law proportional to
$\exp(N\beta\ReTr(QA_e))dQ$ with $A_e$ a sum of six unitary matrices that
depend on $Q'$. For $g\in G$,
\[
|\Tr(Q(g-1)A_e)|\le\|g-1\|\,\Tr|A_eQ|\le6N\|g-1\|,
\]
so the exponent satisfies the hypothesis of Lemma~\ref{lem:onelinkdecay} with
$\Lambda=6N^2\beta$. The function $\Phi(\pi(Q_{e_1})\pi(Q_{e_2})\cdots
\pi(Q_{e_n});Q')$ is linear in each $\pi(Q_e)$, $e\in\{e_2,\dots,e_{n-1}\}$,
separately, so its conditional expectation given $Q'$ is
$\Phi(\pi(Q_{e_1})M_{e_2}\cdots M_{e_{n-1}}\pi(Q_{e_n});Q')$ with
$M_e=\int\pi\,d\nu_e$, and its modulus is at most
$C\prod_e\|M_e\|$. The lemma bounds each factor.
\end{proof}

\subsection{The limit of the transporters}

Let $a_k=2^{-k}a_0$ and let $\beta_k\ge0$ be couplings with
\begin{equation}\label{eq:tuningT}
a_k\max(1,\beta_k)^2\longrightarrow0.
\end{equation}
This holds whenever $\beta_k$ grows at most like a power of $\log(1/a_k)$, in
particular for the tuning $\beta_k\sim\log(1/a_k)$ suggested by asymptotic
freedom in four dimensions. Fix $j$ and a finite set $B$ of bonds of the
lattice $a_j\Z^4$. For $k\ge j+2$ a bond $b\in B$ is a straight segment of
$2^{k-j}$ links of $a_k\Z^4$; let $U_b^{(k)}$ be the ordered product of the
link variables along it, and let $\mu^{(k)}$ be the Wilson measure at
$\beta_k$ on a torus of any side large enough to contain $B$.

\begin{theorem}\label{thm:haarlimit}
Under~\eqref{eq:tuningT}, for every continuous $F\colon G^B\to\C$,
\[
\mu^{(k)}\bigl(F\bigl((U_b^{(k)})_{b\in B}\bigr)\bigr)\longrightarrow
\int F\prod_{b\in B}dQ_b\qquad(k\to\infty).
\]
The parallel transporters along the bonds of any fixed lattice become
independent and Haar distributed.
\end{theorem}

\begin{proof}
By the Peter--Weyl theorem the finite linear combinations of the products
of matrix coefficients $\prod_b\ip{u_b}{\pi_b(Q_b)v_b}$, with $\pi_b$ irreducible and $u_b,v_b$ unit
vectors, are uniformly dense in $C(G^B)$, and both sides are continuous for
the uniform norm; it is enough to treat one product. If every $\pi_b$ is
trivial the two sides are the same constant. Otherwise let $\pi_{b_0}$ be
nontrivial; the Haar integral vanishes. Distinct bonds of $a_j\Z^4$ share no
link of $a_k\Z^4$, so Proposition~\ref{prop:segments} applies to the segment
$b_0$ with $n=2^{k-j}$,
\[
\Phi(A;Q')=\ip{u_{b_0}}{Av_{b_0}}\prod_{b\ne b_0}\ip{u_b}{\pi_b(U_b^{(k)})v_b},
\qquad C=1,
\]
and bounds the left side by
\[
\Bigl(1-\frac{c}{\max(1,6N^2\beta_k)^2}\Bigr)^{2^{k-j}-2}
\le\exp\Bigl(-\frac{c\,(a_j/a_k-2)}{\max(1,6N^2\beta_k)^2}\Bigr),
\qquad c=c_{\pi_{b_0}},
\]
which tends to $0$ by~\eqref{eq:tuningT}.
\end{proof}

Gauge invariance of $F$ was not used. For the trace of the holonomy of a
lattice loop the theorem says that every Wilson loop of fixed physical size,
in every nontrivial irreducible representation, tends to its value at
$\beta=0$, which is zero for a simple loop. This is the perimeter divergence
in rigorous form, and it is the reason why the fields of a continuum limit
cannot be the transporters or the loops themselves: in the limit they carry
no correlation at all. Condition~\eqref{eq:tuningT} is where the dimension
enters; in two dimensions the coupling grows like $a^{-2}$ and Wilson loops
have nontrivial limits.

\subsection{Theories constant along a plane}

Two-dimensional Yang--Mills theory is exactly solvable and has a gap, and one
may try to read it as a theory on $\R^4$ that does not depend on two of the
coordinates.

\begin{proposition}\label{prop:planar}
Let $(\mathcal H,U,\Omega,\phi)$ be a Wightman theory on $\R^{1+3}$ with a
unique vacuum that is cyclic for the fields, and suppose $\phi(x+c)=\phi(x)$
for every $c$ in a spacelike two-plane $E$. Then $\mathcal H=\C\Omega$.
\end{proposition}

\begin{proof}
For $c\in E$ the unitary $U(c)$ commutes with the fields and fixes $\Omega$,
so $U(c)=1$ by cyclicity. The joint spectrum $\Sigma$ of the energy and
momentum is then contained in the two-plane $E^\perp$, which in suitable
coordinates is $\{p_2=p_3=0\}$. It is also a closed Lorentz-invariant subset
of the closed forward cone. If $p\in\Sigma$ and $p\ne0$, then either
$p_1\ne0$, and a rotation in the plane of the first and second coordinates
moves $p$ out of $\{p_2=0\}$, or $p_1=0$ and $p_0>0$, and a boost in the
second direction does. Hence $\Sigma=\{0\}$, every vector is translation
invariant, and the uniqueness of the vacuum gives $\mathcal H=\C\Omega$.
\end{proof}

Independent copies of the two-dimensional theory on the parallel planes
$E^\perp+z$ fare no better: the average $\epsilon^2\sum_zg(\epsilon z)X_z$ of
independent identically distributed variables against a test function has
variance $\epsilon^4\sum_zg(\epsilon z)^2\Var X\to0$, so every field smeared
across the planes is a constant, and no rotation mixes the planes.

\subsection{Fixed lattice spacing}

The lattice theory of Theorem~\ref{thm:A} has a gap at a fixed spacing $a$,
and its fields may be regarded as distributions on $\R^4$ supported on
$a\Z^4$. The Schwinger functions are then invariant under the lattice
translations and the hyperoctahedral group only. Averaging over a uniform
shift restores translation invariance, but not rotation invariance, and it
destroys reflection positivity. The mechanism is visible in one dimension.
Let $(\xi_n)_{n\in\Z}$ be the stationary Gaussian sequence with covariance
$\lambda^{|n-m|}$, $0<\lambda<1$, which is reflection positive, and let
$\phi_u=\sum_n\xi_n\delta_{n+u}$ with $u$ uniform on $[0,1)$. The averaged
two-point function is $\sum_{k\in\Z}\lambda^{|k|}\delta(t-s-k)$, and for a
test function $f$ supported in $(0,1)$ the reflection form is
\[
\sum_{k}\lambda^{|k|}\int\overline{f(t)}f(-t-k)\,dt
=\lambda\int_0^1\overline{f(t)}f(1-t)\,dt,
\]
only $k=-1$ contributing. For $f$ odd about $\frac12$ this equals
$-\lambda\int|f|^2<0$.

\medskip
The three shortcuts fail for one reason. The mass gap is a statement about
the renormalised local fields at a distance fixed in physical units, and a
construction that avoids renormalisation, or avoids four dimensions, or
avoids the continuum, does not produce those fields. What is left open in
this paper is the statement of Theorem~F: the Wightman limit~(W) and the
one-scale bound~\textup{(B$^{*}$)} of Section~\ref{ssec:onehyp}.
